\documentclass{rgtmg} % style de\usepackage{savesym}
\usepackage[utf8]{inputenc}
\usepackage{t1enc}
\usepackage{dfadobe}
\usepackage{cite}
\usepackage[french]{babel}

\title{Réévalutation de modèle basée sur les règles de transformation de graphe}

\author[M. Gaide et al.]{%
	\parbox{\textwidth}{\centering M. Gaide$^{1}$, D. Marcheix$^{2}$, A.
		Arnould$^{3}$, X. Skapin$^{3}$, H. Belhaouari$^{3}$, and S. Jean$^{2}$
		\\
		\vspace{1em}
		$^1$ISAE-ENSMA Poitiers, Université de Poitiers, LIAS, Poitiers, France\\
		$^2$Université de Poitiers, ISAE-ENSMA Poitiers, LIAS, Poitiers, France\\
		$^3$Université de Poitiers, Univ. Limoges, CNRS, XLIM, Poitiers, France}}

\begin{document}
\maketitle

\begin{abstract}
  Dans ce papier, nous étendons l'étude du problème de nommage aux systèmes de
  modélisations basés sur les règles de transformation de graphe. %
  Pour cela, nous proposons une méthode de nommage persistant s'appuyant sur les
  cartes généralisées et la formalisation des opérations par des règles de
  transformation de graphe. %
  Ces formalismes permettant la caractérisation unique et homogène des entités
  topologiques en toutes dimensions. %

  Afin de réévaluer des objets, la plupart des méthodes existantes nécessitent
  de suivre la totalité, ou presque, des entités topologiques qui les
  composent. %
  Aussi, ces mêmes méthodes ne considèrent le problème du nommage que du point
  de vue de la modification des paramètres dans une spécification
  paramétriques. %

  Avec notre méthode, non seulement nous abordons le nommage persistant dans le
  cadre habituel de l'édition de paramètre, mais nous prenons aussi en compte
  l'édition de la spécification paramétrique elle-même (ajout, suppression et
  déplacement d'opération). %
  De plus, nous utilisons des graphes orientés acycliques pour représenter
  l'histoire des entités passées en paramètres dans les opérations jusqu'à leurs
  origines. %

  Une telle représentation permet la mise en place d'un système de nommage
  persistant efficace prenant en compte les différents impacts dû aux
  modifications d'une spécification paramétrique. %
\end{abstract}

\keywords{Modélisation géométrique à base topologique; Règles de transformation
  de graphe; Nommage Persistant; Réévaluation; Cartes généralisées;}

\end{document}
